\section{Complementos sobre las extensiones de álgebras}
\label{section:0.18-es}

Este parágrafo reúne cierto número de construcciones funtoriales sobre
los anillos, que se utilizarán repetidamente en los \S\S~19 y 20~; no
contiene ningún resultado no trivial.

\subsection{Imágenes inversas de anillos aumentados}
\label{subsection:0.18.1-es}

\begin{env}[18.1.1]
\label{0.18.1.1-es}
Dado un anillo $A$ (no necesariamente conmutativo), la categoría de los
\emph{$A$-anillos} tiene por objetos los pares $(B,\rho)$ formados por un anillo
$B$ y un homomorfismo de anillos $\rho:A\to B$, y por morfismos (llamados también
\emph{$A$-homomorfismos}) los homomorfismos de anillos $u:B\to C$ tales que,
si $\rho:A\to B$ y $\sigma:A\to C$ son los homomorfismos (llamados
\emph{estructurales}) que definen la estructura de $A$-anillo sobre $B$ y $C$
respectivamente, el diagrama
\begin{equation}
\label{0.18.1.1.1-es}
\tag{18.1.1.1}
  \xymatrix{
    B\ar[r]^u & C\\
    & A\ar[lu]^\rho\ar[u]_\sigma
  }
\end{equation}
sea conmutativo. El \emph{núcleo} $\mathfrak{J}$ de $u$ es un ideal bilateral
que es un $A$-\emph{bimódulo} (para $\rho$).

Si $f:A'\to A$ es un homomorfismo de anillos, para todo $A$-anillo
$(B,\rho)$, $(B,\rho\circ f)$ es un $A'$-anillo, y si $u:B\to C$ es un
$A$-homomorfismo del $A$-anillo $(B,\rho)$
\oldpage[0\textsubscript{IV-1}]{52}
en el $A$-anillo $(C,\sigma)$, es también un $A'$-homomorfismo del
$A'$-anillo $(B,\rho\circ f)$ en el $A'$-anillo $(C,\sigma\circ f)$~; se
define así un funtor canónico de la categoría de los $A$-anillos en la de
los $A'$-anillos.
\end{env}

\begin{env}[18.1.2]
\label{0.18.1.2-es}
Sean $B$, $E$, $F$ tres $A$-anillos, $f:E\to B$, $g:F\to B$ dos
$A$-homomorfismos. Recordemos que se llama \emph{producto fibrado} de $E$ y
$F$ sobre $B$ (para los $A$-homomorfismos $f$ y $g$) y se denota por
$E\times_BF$ el subanillo $G$ del anillo producto $E\times F$, formado por los
pares $(x,y)$ tales que $f(x)=g(y)$~; las restricciones $p_1:G\to E$,
$p_2:G\to F$ de las proyecciones $\operatorname{pr}_1$ y
$\operatorname{pr}_2$ a $G$ siguen llamándose las \emph{proyecciones
canónicas}~; la estructura de $A$-anillo de $G$ está definida por
el homomorfismo $\alpha\mapsto(\rho(\alpha),\sigma(\alpha))$ (donde
$\rho:A\to E$ y $\sigma:A\to F$ son los homomorfismos estructurales), que
envía efectivamente $A$ a $G$ en virtud de (18.1.1). La propiedad
característica de $E\times_BF$ es que, para todo par de
$A$-homomorfismos $u:C\to E$, $v:C\to F$ tales que
$f\circ u=g\circ v$, existe un único $A$-homomorfismo $w:C\to G$
tal que $u=p_1\circ w$, $v=p_2\circ w$. También puede decirse que
$E\times_BF$ es el \emph{límite proyectivo} del sistema proyectivo formado por $B$,
$E$, $F$ y los $A$-homomorfismos $f$, $g$, en la categoría de los
$A$-anillos ($0_{\mathrm{III}}$, 8.1.9).
\end{env}

\begin{env}[18.1.3]
\label{0.18.1.3-es}
Sea $\mathfrak{J}$ el ideal bilateral de $E$, núcleo de $f$~; es inmediato
que el núcleo $\mathfrak{J}'$ de $p_2:G\to F$ es el ideal formado por los elementos
$(x,0)$ donde $x\in\mathfrak{J}$~; la restricción
$i_1:\mathfrak{J}'\to\mathfrak{J}$ de $p_1$ es, por tanto, un
\emph{isomorfismo} de anillo \emph{sin elemento unidad}, y también un
isomorfismo de $A$-bimódulos. Del mismo modo, si $\mathfrak{R}$ es el núcleo de
$g$, el núcleo $\mathfrak{R}'$ de $p_1:G\to E$ es el ideal de los elementos
$(0,y)$ donde $y\in\mathfrak{R}$, y la restricción
$i_2:\mathfrak{R}'\to\mathfrak{R}$ de $p_2$ es un isomorfismo (en el mismo
sentido). Finalmente, es claro que el núcleo del $A$-homomorfismo
$q=f\circ p_1=g\circ p_2$ de $E\times_BF$ en $B$ es
$\mathfrak{J}\times\mathfrak{R}=\mathfrak{J}'\oplus\mathfrak{R}'$, de modo
que se tiene el diagrama conmutativo
\begin{equation}
\label{0.18.1.3.1-es}
\tag{18.1.3.1}
  \xymatrix@C=2.2pc@R=1.8pc{
    0\ar[dr] & & 0\ar[d] & 0\ar[d]\\
    & \mathfrak{J}'\oplus\mathfrak{R}'\ar[r]\ar[d]\ar[dr]
      & \mathfrak{R}'\ar[r]^{i_2}_{\sim}\ar[d]^j
      & \mathfrak{R}\ar[d]^{p_2}\\
    0\ar[r] & \mathfrak{J}'\ar[r]\ar[d]^{i_1}_{\sim}
      & G\ar[r]\ar[d]^{p_1}\ar[dr]^q & F\ar[d]^g\\
    0\ar[r] & \mathfrak{J}\ar[r] & E\ar[r]_f & B
  }
\end{equation}

Las definiciones y los resultados de (18.1.2) y (18.1.3) se extienden inmediatamente al
producto fibrado de una \emph{familia} arbitraria $(E_\lambda)_{\lambda\in I}$ de
$A$-anillos definida por una familia de $A$-homomorfismos
$f_\lambda:E_\lambda\to B$. Dejamos al lector la formulación de estos resultados.
\end{env}

\begin{env}[18.1.4]
\label{0.18.1.4-es}
De acuerdo con la terminología de (M, VIII), llamaremos
\emph{$A$-anillo aumentado sobre $B$} a un $A$-anillo $E$ provisto de un
$A$-homomorfismo \emph{sobreyectivo} (llamado \emph{aumentación} de $E$),
$f:E\to B$~; el núcleo $\mathfrak{J}$ de $f$ se llama
\emph{ideal de aumentación}. Se dice que el $A$-anillo aumentado $E$ es
\emph{trivial} si existe un $A$-homomorfismo de $A$-anillos $s:B\to E$
que es \emph{inverso por la derecha} de la aumentación $f:E\to B$ (dicho de otro modo,
$f\circ s=1_B$). La sucesión exacta de $A$-bimódulos
\[
  0\longrightarrow\mathfrak{J}\longrightarrow E
  \xrightarrow{f}B\longrightarrow0
\]
\oldpage[0\textsubscript{IV-1}]{53}
es entonces \emph{escindida}~; dicho de otro modo, se puede identificar el $A$-bimódulo
$E$ con $B\times\mathfrak{J}$, y la multiplicación en $E$ viene dada entonces
por
\[
  (b,z)(b',z')=(bb',bz'+zb'+zz'),
\]
considerándose $\mathfrak{J}$ como $B$-bimódulo por medio de $s:B\to E$.
\end{env}

\begin{env}[18.1.5]
\label{0.18.1.5-es}
Con las notaciones de (18.1.2), supongamos que el $A$-homomorfismo $f$ es
\emph{sobreyectivo}, dicho de otro modo, que hace de $E$ un $A$-anillo \emph{aumentado sobre
$B$}. Entonces es claro que $p_2:G\to F$ es también \emph{sobreyectivo},
dicho de otro modo, define sobre $G$ una estructura de $A$-anillo \emph{aumentado sobre
$F$}, que se llama \emph{imagen inversa} por $g:F\to B$ del anillo
aumentado $E$.
\end{env}

\begin{proposition}[18.1.6]
\label{0.18.1.6-es}
Para que la imagen inversa $G=E\times_BF$ por $g:F\to B$ del $A$-anillo
aumentado $E$ sea un $A$-anillo aumentado trivial, es necesario y suficiente que
exista un $A$-homomorfismo $u:F\to E$ que haga conmutativo el diagrama
\[
  \xymatrix{
    & F\ar[ld]_u\ar[d]^g\\
    E\ar[r]_f & B
  }
\]
\end{proposition}

\begin{proof}
La condición es evidentemente necesaria, tomando $u=p_1\circ s$, donde
$s:F\to G$ es un $A$-homomorfismo inverso por la derecha de $p_2$. Recíprocamente,
si existe un $A$-homomorfismo $u$ que verifica la condición del enunciado,
la existencia del inverso por la derecha $s$ de $p_2$ resulta de la propiedad
universal del producto fibrado (18.1.2) aplicada a los $A$-homomorfismos
$u:F\to E$ y $1_F:F\to F$.

Este resultado implica, en particular, que si $E$ es un anillo aumentado
\emph{trivial}, ocurre lo mismo con todas sus imágenes inversas.
\end{proof}

\begin{env}[18.1.7]
\label{0.18.1.7-es}
Retomemos la situación descrita en (18.1.2) y (18.1.3)~; se tiene evidentemente sobre
$\mathfrak{R}$ una estructura de $G$-bimódulo, procedente de su estructura de
$F$-bimódulo y del homomorfismo de anillos $p_2:G\to F$. Como $i_2$ es
biyectivo, se tiene además una inyección
$\theta=j\circ i_2^{-1}:\mathfrak{R}\to G$, que es un homomorfismo de
$G$-bimódulos. Vamos a ver, recíprocamente, que, cuando se supone además
que $f$ y $g$ son \emph{sobreyectivos}, dicho de otro modo, que $E$ y $F$ son
$A$-anillos \emph{aumentados} sobre $B$, el dato de tal homomorfismo
$\theta$ permite reconstruir el anillo aumentado $E$ (sobre $B$) a partir del
anillo aumentado $G$ (sobre $F$).

Más precisamente, démonos un $A$-anillo $F$ aumentado sobre $B$ y un
$A$-anillo $G$ aumentado sobre $F$, y denotemos por $\mathfrak{R}$ y
$\mathfrak{J}'$, respectivamente, sus ideales de aumentación~:
\begin{equation}
\label{0.18.1.7.1-es}
\tag{18.1.7.1}
  \xymatrix@C=2pc@R=1.5pc{
    & & 0\ar[d] & & \\
    & & \mathfrak{R}\ar[ld]^\theta\ar[d] & & \\
    0\ar[r] & \mathfrak{J}'\ar[r] & G\ar[r]_h & F\ar[r]\ar[d]^g & 0\\
    & & & B\ar[d] & \\
    & & & 0 &
  }
\end{equation}

\oldpage[0\textsubscript{IV-1}]{54}
El homomorfismo $h$ define sobre todos los ideales de $F$ (y, en particular, sobre
$\mathfrak{R}$) una estructura de $G$-bimódulo. Sea
$\theta:\mathfrak{R}\to G$ un homomorfismo de $G$-bimódulos que haga
conmutativo el diagrama (18.1.7.1)~; esto implica que $\theta$ es
\emph{inyectivo} y que $\mathfrak{J}'\cap\theta(\mathfrak{R})=0$~; como
$\theta(\mathfrak{R})$ es un ideal bilateral de $G$, y
$h(\theta(\mathfrak{R}))=\mathfrak{R}$, el homomorfismo compuesto
$g\circ h:G\to B$ se factoriza como
$G\xrightarrow{q}G/\theta(\mathfrak{R})\xrightarrow{f}B$, donde $f$ es
sobreyectivo~; además, la imagen $\mathfrak{J}$ de $\mathfrak{J}'$ en
$E=G/\theta(\mathfrak{R})$ es el núcleo de $f$ (siendo el de $g\circ h$
$\mathfrak{J}'\oplus\theta(\mathfrak{R})$), y la restricción a
$\mathfrak{J}'$ del homomorfismo canónico $q:G\to E$ es inyectiva. Se
puede, por tanto, formar el \emph{producto fibrado} $G'=E\times_BF$, y como los dos
$A$-homomorfismos $q:G\to E$ y $h:G\to F$ son tales que
$f\circ q=g\circ h$, definen un único $A$-homomorfismo $u:G'\to G$
por la propiedad universal de $G'$ (18.1.2)~; vamos a ver que $u$ es
\emph{biyectivo}. Basta demostrarlo cuando $u$ se considera como
homomorfismo de $A$-bimódulos~; pero se observará entonces que $u$ es compatible
con las filtraciones finitas sobre $G'$ y $G$ formadas respectivamente por $G'$
y $\mathfrak{J}''=\operatorname{Ker}(G'\to F)$, y por $G$ y
$\mathfrak{J}'$~; además, se ha visto en (18.1.3) que
$\operatorname{gr}_0u:F\to F$ es la identidad y que
$\operatorname{gr}_1u:\mathfrak{J}''\to\mathfrak{J}'$ es biyectivo, por lo que $u$
mismo es biyectivo (Bourbaki, \emph{Alg. comm.}, cap.~III, \S~2,
n\textsuperscript{o}~8, cor.~3 del teor.~1).
\end{env}

\subsection{Extensiones de un anillo por un bimódulo}
\label{subsection:0.18.2-es}

\begin{env}[18.2.1]
\label{0.18.2.1-es}
Sean $E$ un $A$-anillo aumentado sobre $B$, $f:E\to B$ la aumentación,
$\mathfrak{J}=\operatorname{Ker}(f)$ el ideal de aumentación. Si se tiene
$\mathfrak{J}^2=0$, $\mathfrak{J}$ no es solamente un $E$-bimódulo, sino
también un $B$-\emph{bimódulo}, puesto que $B$ es isomorfo a
$E/\mathfrak{J}$~; más precisamente, todo $b\in B$ es de la forma $f(x)$
con $x\in E$, y si $z,z'$ están en $\mathfrak{J}$ se tiene
$(x+z)z'=xz'$ (resp. $z'(x+z)=z'x$), de modo que el valor de $xz'$ (resp.
$z'x$) no depende del elemento $x$ de $f^{-1}(b)$, y puede escribirse $bz'$
(resp. $z'b$), lo que define la estructura de $B$-bimódulo considerada.
Recíprocamente, si $\mathfrak{J}$ está provisto de una estructura de $B$-bimódulo tal
que $xz'=f(x)z'$ y $z'x=z'f(x)$ para $x\in E$ y
$z'\in\mathfrak{J}$, es claro que $\mathfrak{J}^2=0$.
\end{env}

\begin{env}[18.2.2]
\label{0.18.2.2-es}
Se llama \emph{$A$-extensión de un $A$-anillo $B$ por un $B$-bimódulo $L$}
a una \emph{sucesión exacta} de homomorfismos de $A$-bimódulos
\[
  0\longrightarrow L\xrightarrow{j}E\xrightarrow{f}B\longrightarrow0
\]
donde $E$ es un $A$-anillo, $f$ un $A$-homomorfismo de anillos y se tiene, para
$x\in E$ y $z\in L$,
\[
  j(f(x)z)=xj(z),\qquad j(zf(x))=j(z)x
\]
de donde se deduce, en virtud de (18.2.1), que $j(L)$ es un ideal bilateral de cuadrado nulo de $E$.
Por abuso de lenguaje, se dirá también que $E$ es una extensión de $B$ por $L$.
Se dice que dos $A$-extensiones $E$, $E'$ de $B$ por $L$ son
\emph{$A$-equivalentes} si existe un isomorfismo de $A$-anillos
$u:E\xrightarrow{\sim}E'$ (llamado también \emph{$A$-equivalencia de
$A$-extensiones}) que hace conmutativo el diagrama
\begin{equation}
\label{0.18.2.2.1-es}
\tag{18.2.2.1}
  \xymatrix@C=2.1pc@R=1.2pc{
    & & E\ar[dr]\ar[dd]_u\ar@{}[dd]|{\wr} & & \\
    0\ar[r] & L\ar[ur]\ar[dr] & & B\ar[r] & 0\\
    & & E'\ar[ur] & &
  }
\end{equation}
\end{env}

\oldpage[0\textsubscript{IV-1}]{55}

\begin{env}[18.2.3]
\label{0.18.2.3-es}
Se dice que una $A$-extensión $E$ de $B$ por un $B$-bimódulo $L$ es
\emph{$A$-trivial} si $E$ es un $A$-anillo aumentado (sobre $B$) trivial
(18.1.4). Se define sobre el $A$-bimódulo producto $B\times L$ una estructura de
$A$-anillo poniendo $(x,s)(y,t)=(xy,xt+sy)$, como se comprueba inmediatamente,
y es inmediato que las aplicaciones canónicas $j:L\to B\times L$ y
$f:B\times L\to B$ definen una $A$-extensión de $B$ por $L$ que es
$A$-trivial, siendo la aplicación canónica $g:B\to B\times L$ un
$A$-homomorfismo inverso por la derecha de $f$. Se dice que esta extensión es
\emph{la extensión trivial tipo} de $B$ por $L$ y se denota por $D_B(L)$; es
inmediato que toda $A$-extensión $A$-trivial de $B$ por $L$ es
$A$-equivalente a $D_B(L)$.

Se observará que toda extensión del propio $A$-anillo $A$ por un $A$-bimódulo
es necesariamente $A$-trivial.
\end{env}

\begin{env}[18.2.4]
\label{0.18.2.4-es}
Dadas dos $A$-extensiones
$L\xrightarrow{j}E\xrightarrow{f}B$,
$L'\xrightarrow{j'}E'\xrightarrow{f'}B'$, un \emph{morfismo} de la primera
en la segunda es, por definición, una terna de homomorfismos de
$A$-bimódulos $(u,v,w)$ tal que el diagrama
\begin{equation}
\label{0.18.2.4.1-es}
\tag{18.2.4.1}
  \xymatrix{
    0\ar[r] & L\ar[r]^j\ar[d]_w & E\ar[r]^f\ar[d]_u
      & B\ar[r]\ar[d]_v & 0\\
    0\ar[r] & L'\ar[r]_{j'} & E'\ar[r]_{f'} & B'\ar[r] & 0
  }
\end{equation}
sea conmutativo, siendo $u$ y $v$ $A$-homomorfismos \emph{de anillos} y
$w$ tal que $w(bz)=v(b)w(z)$ y $w(zb)=w(z)v(b)$ para $z\in L$ y
$b\in B$ (en otros términos, el par $(v,w)$ constituye un
\emph{dihomomorfismo} del $B$-bimódulo $L$ en el $B'$-bimódulo $L'$); es
claro que, si $(u',v',w')$ es un morfismo de
$L'\to E'\to B'$ en una $A$-extensión $L''\to E''\to B''$,
$(u'\circ u,v'\circ v,w'\circ w)$ es también un morfismo, lo que justifica la
terminología.

La consideración de los dos cuadrados conmutativos del diagrama (18.2.4.1) va a
conducirnos a dos operaciones sobre las extensiones de $A$-anillos.
\end{env}

\begin{env}[18.2.5]
\label{0.18.2.5-es}
En primer lugar, consideremos una $A$-extensión $E'$ de $B'$ por $L'$
\[
  0\longrightarrow L'\xrightarrow{j'}E'\xrightarrow{f'}B'\longrightarrow0
\]
y un $A$-homomorfismo de anillos $v:B\to B'$, y sea
$F=E'\times_{B'}B$ \emph{la imagen inversa} por $v$ del $A$-anillo aumentado
$E'$ (18.1.5), de modo que se tiene un diagrama conmutativo
\[
  \xymatrix{
    0\ar[r] & L_0\ar[r]\ar[d]^i\ar@{}[d]|{\wr} & F\ar[r]^{p_2}\ar[d]^{p_1}
      & B\ar[r]\ar[d]^v & 0\\
    0\ar[r] & L'\ar[r]_{j'} & E'\ar[r]_{f'} & B'\ar[r] & 0
  }
\]
cuyas filas son exactas, siendo $p_1$ y $p_2$ los homomorfismos
canónicos; se ha visto (18.1.3) que $i$ es biyectivo y de la
definición (18.1.2) resulta también que $L_0^2=0$, de modo que se puede considerar que $F$
es una $A$-extensión de $B$ por $L'$, que se denomina \emph{imagen
inversa por $v$} de la extensión $E'$ de $B'$ por $L'$ (considerándose
naturalmente $L'$ como $B$-bimódulo mediante el homomorfismo
de anillos $v$). El carácter funtorial
\oldpage[0\textsubscript{IV-1}]{56}
del producto fibrado respecto de cada uno de los factores muestra además que, si se
tiene un morfismo de dos extensiones de $B'$
\[
  \xymatrix{
    0\ar[r] & L'_1\ar[r]\ar[d]^h & E'_1\ar[r]\ar[d]^g
      & B'\ar[r]\ar[d]^{1_{B'}} & 0\\
    0\ar[r] & L'_2\ar[r] & E'_2\ar[r] & B'\ar[r] & 0
  }
\]
se deduce de él un morfismo \emph{imagen inversa por $v$}
\[
  \xymatrix{
    0\ar[r] & L'_1\ar[r]\ar[d]^h & E'_1\times_{B'}B\ar[r]
      \ar[d]^{g\times_{B'}1_B} & B\ar[r]\ar[d]^{1_B} & 0\\
    0\ar[r] & L'_2\ar[r] & E'_2\times_{B'}B\ar[r] & B\ar[r] & 0.
  }
\]
En particular, si $E'_1$ y $E'_2$ son $A$-extensiones
$A$-equivalentes de $B'$ por $L'$, sus imágenes inversas por $v$ son
$A$-extensiones $A$-equivalentes de $B$ por $L'$.

La definición del producto fibrado muestra que, cuando se tiene un morfismo
(18.2.4.1) de $A$-extensiones, este se \emph{factoriza a través de} la imagen
inversa de $E'$ por $v$; más precisamente, existe un
$A$-homomorfismo único $u_0:E\to F=E'\times_{B'}B$ que hace conmutativo el
diagrama
\[
  \xymatrix{
    0\ar[r] & L\ar[r]^j\ar[d]^{w_0} & E\ar[r]^f\ar[d]^{u_0}
      & B\ar[r]\ar[d]^{1_B} & 0\\
    0\ar[r] & L_0\ar[r]\ar[d]^i\ar@{}[d]|{\wr} & F\ar[r]^{p_2}\ar[d]^{p_1}
      & B\ar[r]\ar[d]^v & 0\\
    0\ar[r] & L'\ar[r]_{j'} & E'\ar[r]_{f'} & B'\ar[r] & 0
  }
\]
donde $w_0$ es la restricción de $u_0$ a $L$ y
$p_1\circ u_0=u$, $i\circ w_0=w$.
\end{env}

\begin{env}[18.2.6]
\label{0.18.2.6-es}
Estudiemos en particular las imágenes inversas de extensiones por
$A$-homomorfismos \emph{sobreyectivos}. Consideremos un $A$-homomorfismo
sobreyectivo $v:B\to B'$, una $A$-extensión $F$ de $B$ por un $B$-bimódulo $L$,
y el ideal bilateral $\mathfrak{R}$ de $B$, núcleo de $v$ (que puede ser
considerado como $F$-bimódulo mediante la aumentación $F\to B$); se ha visto
(18.1.7) que todo homomorfismo de $F$-bimódulos
$\theta:\mathfrak{R}\to F$ que hace conmutativo el diagrama
\[
  \xymatrix{
    & \mathfrak{R}\ar[ld]^\theta\ar[d]\\
    F\ar[r] & B
  }
\]
determina una extensión $E'$ de $B'=B/\mathfrak{R}$ por $L$ cuya imagen
inversa por $v:B\to B'$ es equivalente a $F$, y que toda $A$-extensión
$E'$ de $B'$ por $L$ que tenga esta última propiedad se obtiene así (salvo
$A$-equivalencia). Además, para que dos homomorfismos
$\theta_1,\theta_2$ de $F$-bimódulos de $\mathfrak{R}$ en $F$ den dos
  $A$-extensiones $A$-equivalentes de $B$ por $L$, es necesario y suficiente que
exista una $A$-equivalencia $u$ de la $A$-extensión $F$ sobre sí misma tal que
\oldpage[0\textsubscript{IV-1}]{57}
$\theta_2=u\circ\theta_1$; esto resulta inmediatamente de lo que se ha visto en
(18.2.5) y de la definición de la biyección canónica de $F$ sobre el producto
fibrado $E'\times_{B'}B$ (18.1.7).
\end{env}

\begin{env}[18.2.7]
\label{0.18.2.7-es}
Consideremos ahora el cuadrado izquierdo de (18.2.4.1), y recordemos primero
la noción de \emph{suma amalgamada} en la categoría de los $A$-bimódulos:
dados tres $A$-bimódulos $X$, $Y$, $Z$ y dos $A$-homomorfismos
$f:X\to Y$, $g:X\to Z$, la suma amalgamada $Y\oplus_XZ$ es el \emph{límite
inductivo} del sistema inductivo formado por $X$, $Y$, $Z$ y los
$A$-homomorfismos $f$, $g$ en la categoría de los $A$-bimódulos
($0_{\mathrm{III}}$, 8.1.11). Se define este $A$-bimódulo como cociente del
producto $Y\times Z$ por el sub-$A$-bimódulo $M$ imagen de $X$ por
el homomorfismo $x\mapsto(f(x),-g(x))$. Su propiedad característica es que,
para todo par de homomorfismos de $A$-bimódulos $u:Y\to T$, $v:Z\to T$
tales que $u\circ f=v\circ g$, existe un único homomorfismo
$w:Y\oplus_XZ\to T$ tal que $u=w\circ j_1$ y $v=w\circ j_2$, donde
$j_1:Y\to Y\oplus_XZ$ y $j_2:Z\to Y\oplus_XZ$ son las aplicaciones
canónicas.
\end{env}

\begin{env}[18.2.8]
\label{0.18.2.8-es}
Consideremos ahora una $A$-extensión $E$ de $B$ por un $B$-bimódulo
$L$:
\[
  0\longrightarrow L\xrightarrow{j}E\xrightarrow{f}B\longrightarrow0
\]
y sea, por otra parte, $w:L\to L'$ un homomorfismo de $B$-bimódulos. Sea $H$
el $A$-bimódulo suma amalgamada $E\oplus_LL'$; mostremos cómo se puede dotar
este $A$-bimódulo de una estructura de $A$-anillo y definir una $A$-extensión
\[
  0\longrightarrow L'\longrightarrow H\longrightarrow B\longrightarrow0.
\]

Observemos para ello que $L'$ está dotado de una estructura de $E$-\emph{bimódulo} por
medio del homomorfismo de aumentación $E\to B$; por tanto, se puede formar la
$A$-extensión trivial tipo $G=D_E(L')$ (18.2.3). Consideremos entonces
la aplicación $\theta:z\mapsto(j(z),-w(z))$ de $L$ en $G$; es un
homomorfismo de $G$-bimódulos (considerándose $L$ como $G$-bimódulo por
medio del homomorfismo canónico $\rho:G\to E$). En efecto, para
$(x,z')\in G$ y $z\in L$, se tiene
$(j(z),-w(z))(x,z')=(j(z)x,j(z)z'-w(z)f(x))$; ahora bien,
$j(z)z'=f(j(z))z'=0$ por definición de la estructura de $E$-bimódulo sobre
$L'$, y $w(z)f(x)=w(zf(x))$ por definición de la estructura de $B$-bimódulo
sobre $L'$; se comprueba del mismo modo que $\theta$ es un homomorfismo de
$G$-módulos izquierdos. Se puede aplicar entonces al diagrama conmutativo
\[
  \xymatrix@C=2pc@R=1.4pc{
    & & 0\ar[d] & & \\
    & & L\ar[ld]^\theta\ar[d]^j & & \\
    0\ar[r] & L'\ar[r] & G\ar[r]_p & E\ar[r]\ar[d]^f & 0\\
    & & & B\ar[d] & \\
    & & & 0 &
  }
\]
el resultado de (18.1.7). Como $H=G/\theta(L)$ por definición, nuestra
afirmación es una consecuencia inmediata de (18.1.7).

\oldpage[0\textsubscript{IV-1}]{58}

Se dice que la $A$-extensión $H$ de $B$ por $L'$ es \emph{deducida de $E$ por
medio del homomorfismo} $w:L\to L'$. El carácter funtorial de la suma
amalgamada en cada uno de los sumandos muestra además que, si se tiene un morfismo
de extensiones
\[
  \xymatrix{
    0\ar[r] & L\ar[r]\ar[d]^{1_L}\ar@{}[d]|{\wr}
      & E_1\ar[r]\ar[d]^g & B_1\ar[r]\ar[d]^h & 0\\
    0\ar[r] & L\ar[r] & E_2\ar[r] & B_2\ar[r] & 0
  }
\]
se deduce canónicamente un morfismo de extensiones
\[
  \xymatrix{
    0\ar[r] & L'\ar[r]\ar[d]^{1_{L'}}\ar@{}[d]|{\wr}
      & E_1\oplus_LL'\ar[r]\ar[d]^{g\oplus_L1_{L'}}
      & B_1\ar[r]\ar[d]^h & 0\\
    0\ar[r] & L'\ar[r] & E_2\oplus_LL'\ar[r] & B_2\ar[r] & 0.
  }
\]
En particular, si $E_1$ y $E_2$ son $A$-extensiones $A$-equivalentes de
$B$ por $L$, las extensiones de $B$ por $L'$ que se deducen de ellas mediante $w$
son $A$-equivalentes.

Cuando se tiene un morfismo (18.2.4.1) de $A$-extensiones, este se \emph{factoriza
a través de} la $A$-extensión $H$ de $B$ por $L'$ deducida de $E$ por medio
del homomorfismo $w:L\to L'$ (considerándose $L'$ como $B$-bimódulo mediante
el homomorfismo $v:B\to B'$): en efecto, la definición de la suma amalgamada
muestra que existe un $A$-homomorfismo único $u_0:H\to E'$ de
$A$-bimódulos que hace conmutativo el diagrama
\[
  \xymatrix{
    0\ar[r] & L\ar[r]^j\ar[d]^{w_0} & E\ar[r]^f\ar[d]^{j_1}
      & B\ar[r]\ar[d]^{1_B}\ar@{}[d]|{\wr} & 0\\
    0\ar[r] & L'_0\ar[r]^{j_2}\ar[d]^{j_0}\ar@{}[d]|{\wr}
      & H\ar[r]^{f_0}\ar[d]^{u_0} & B\ar[r]\ar[d]^v & 0\\
    0\ar[r] & L'\ar[r]_{j'} & E'\ar[r]_{f'} & B'\ar[r] & 0
  }
\]
con $u=u_0\circ j_1$, $w=j_0\circ w_0$, siendo $j_1$ y $j_2$ los
homomorfismos canónicos; se comprueba inmediatamente que $u_0$ es también un
homomorfismo de anillos.

Señalemos, por último, las propiedades funtoriales relativas a las extensiones
triviales:
\end{env}

\begin{proposition}[18.2.9]
\label{0.18.2.9-es}
Sean $B$, $B'$ dos $A$-anillos, $L$ un $B$-bimódulo, $L'$ un
$B'$-bimódulo, $v:B\to B'$ un $A$-homomorfismo de anillos, $w:L\to L'$ un
homomorfismo de $A$-bimódulos tal que $(v,w)$ sea un dihomomorfismo de
bimódulos. Entonces existe un $A$-homomorfismo único de anillos
$u:D_B(L)\to D_{B'}(L')$ que hace conmutativos los diagramas
\[
  \xymatrix@C=2.2pc@R=1.6pc{
    D_B(L)\ar[r]^u & D_{B'}(L') & & D_B(L)\ar[r]^u & D_{B'}(L')\\
    B\ar[u]\ar[r]_v & B'\ar[u] & & L\ar[u]\ar[r]_w & L'\ar[u]
  }
\]
donde las flechas verticales son las inyecciones canónicas.
\end{proposition}

\oldpage[0\textsubscript{IV-1}]{59}

\begin{proof}
En efecto, $u$ no puede ser sino la aplicación
$(x,s)\mapsto(v(x),w(s))$, y queda comprobar que es un
$A$-homomorfismo de anillos, lo que resulta trivialmente de la definición
(18.2.3). Se observará que $u$ hace también conmutativo el diagrama
\[
  \xymatrix@C=2.2pc@R=1.5pc{
    D_B(L)\ar[r]^u\ar[d] & D_{B'}(L')\ar[d]\\
    B\ar[r]_v & B'
  }
\]
donde esta vez las flechas verticales son las aumentaciones.
\end{proof}

\begin{proposition}[18.2.10]
\label{0.18.2.10-es}
Sean $B$ un $A$-anillo, $L$ un $B$-bimódulo y $E$ una $A$-extensión de $B$
por $L$. Se define una aplicación biyectiva del conjunto $G$ de los
homomorfismos de $B$-anillos de $B$ en $E$ (es decir, el conjunto de los
$A$-homomorfismos inversos por la derecha de la aumentación $E\to B$) sobre el
conjunto $G'$ de las $A$-equivalencias de $D_B(L)$ sobre $E$, haciendo
corresponder a todo $g\in G$ la $A$-equivalencia
$g':(x,s)\mapsto g(x)+s$; la aplicación inversa hace corresponder a todo
$g'\in G'$ el $B$-homomorfismo $x\mapsto g'(x,0)$.
\end{proposition}

\begin{proof}
Esto resulta inmediatamente de las definiciones.
\end{proof}

\subsection{El grupo de las clases de $A$-extensiones}
\label{subsection:0.18.3-es}

\begin{env}[18.3.1]
\label{0.18.3.1-es}
Consideremos un $A$-anillo $B$ y un $B$-bimódulo $L$ fijados; entonces la
relación «$E$ y $E'$ son $A$-equivalentes» entre las $A$-extensiones $E$, $E'$
de $B$ por $L$ es una relación de equivalencia, y para esta relación se puede
hablar del \emph{conjunto} de las clases de $A$-extensiones $A$-equivalentes de
$B$ por $L$. Para verlo, basta observar que, si para todo $x\in B$, $c_x$ es
un elemento de $E$ cuya imagen en $B$ es $x$, todo $z\in E$ se escribe de una
manera y solo una en la forma $c_x+t$, donde $t\in L$, y se puede escribir
\[
  c_{x+y}=c_x+c_y+\varphi(x,y),\qquad
  c_{xy}=c_xc_y+\psi(x,y),
\]
donde $\varphi(x,y)$ y $\psi(x,y)$ son elementos de $L$, y las aplicaciones
$\varphi$ y $\psi$ de $B\times B$ en $L$ deben verificar condiciones que
expresan que $E$ es un $A$-anillo y que no es necesario escribir aquí. Toda
$A$-extensión de $B$ por $L$ es, por tanto, $A$-equivalente a una
$A$-extensión cuyo conjunto subyacente es $B\times L$, de donde se obtiene
inmediatamente nuestra conclusión.
\end{env}

\begin{env}[18.3.2]
\label{0.18.3.2-es}
Fijado el $A$-anillo $B$, designemos provisionalmente por $T(L)$ el conjunto
de las clases de $A$-extensiones de $B$ por $L$. Para todo $B$-homomorfismo
$w:L\to L'$ de $B$-bimódulos, se define canónicamente una aplicación
$T(w):T(L)\to T(L')$, haciendo corresponder a la clase de una $A$-extensión
$E$ de $B$ por $L$ la clase de la $A$-extensión $E\oplus_LL'$ deducida de $E$
por $w$, en virtud de (18.2.8). Si $w':L'\to L''$ es un segundo homomorfismo
de $B$-bimódulos, se tiene además
\begin{equation}
\label{0.18.3.2.1-es}
\tag{18.3.2.1}
  T(w'\circ w)=T(w')\circ T(w).
\end{equation}
En efecto, se sabe que existe un isomorfismo canónico de $A$-bimódulos
\[
  E\oplus_LL''\xrightarrow{\sim}(E\oplus_LL')\oplus_{L'}L''
\]
en virtud de las propiedades generales de los límites inductivos (cf., por
ejemplo, ($\mathrm{I}$, 3.3.9)), y es inmediato comprobar que se trata
efectivamente de una $A$-equivalencia de $A$-extensiones, de donde se obtiene

\oldpage[0\textsubscript{IV-1}]{60}

la relación (18.3.2.1). Si $\mathbf C_B$ designa la categoría de los
$B$-bimódulos, se ve que se ha definido así un \emph{funtor covariante}
$T:\mathbf C_B\to\mathbf{Ens}$.
\end{env}

\begin{env}[18.3.3]
\label{0.18.3.3-es}
Consideremos ahora una familia $(L_\alpha)_{\alpha\in I}$ de $B$-bimódulos y
su producto $L=\prod_{\alpha\in I}L_\alpha$; las proyecciones
$\operatorname{pr}_\alpha:L\to L_\alpha$ definen aplicaciones
\[
  T(\operatorname{pr}_\alpha):T(L)\to T(L_\alpha),
\]
de donde se obtiene una aplicación canónica
\begin{equation}
\label{0.18.3.3.1-es}
\tag{18.3.3.1}
  \prod_\alpha T(\operatorname{pr}_\alpha):T(L)
  \longrightarrow\prod_{\alpha\in I}T(L_\alpha).
\end{equation}
Vamos a ver que esta aplicación es \emph{biyectiva}. En efecto, para todo
$\alpha\in I$, sea $E_\alpha$ una $A$-extensión de $B$ por $L_\alpha$, y sea
$F$ el \emph{producto fibrado} de los $E_\alpha$ sobre $B$ (18.1.3); es
inmediato que $F$ es una $A$-extensión de $B$ por
$L=\prod_{\alpha\in I}L_\alpha$ y que, si se reemplaza cada $E_\alpha$ por una
$A$-extensión $A$-equivalente $E'_\alpha$, el producto fibrado $F'$ de los
$E'_\alpha$ sobre $B$ es $A$-equivalente a $F$. Se ha definido así una
aplicación $\prod_{\alpha\in I}T(L_\alpha)\to T(L)$, y es claro que esta
aplicación es la inversa de (18.3.3.1), de donde se sigue nuestra afirmación.
Se identificará canónicamente $T(L)$ con
$\prod_{\alpha\in I}T(L_\alpha)$ mediante la aplicación (18.3.3.1). Se
comprueba además inmediatamente que, si $(L'_\alpha)_{\alpha\in I}$ es una
segunda familia de $B$-bimódulos y, para cada $\alpha$,
$w_\alpha:L_\alpha\to L'_\alpha$ es un homomorfismo de $B$-bimódulos,
entonces, poniendo
$w=\prod_\alpha w_\alpha:L\to L'=\prod_{\alpha\in I}L'_\alpha$, $T(w)$ se
identifica con $\prod_{\alpha\in I}T(w_\alpha)$ cuando se efectúa la
identificación precedente.
\end{env}

\begin{env}[18.3.4]
\label{0.18.3.4-es}
Sentado esto, para un $B$-bimódulo $L$, la adición es un homomorfismo
$s:L\times L\to L$ de $B$-bimódulos, y lo mismo sucede con la simetría
$t:L\to L$ de la ley aditiva de $L$. Se deduce de ello una ley de composición
\[
  T(s):T(L)\times T(L)\longrightarrow T(L)
\]
en $T(L)$, en virtud de (18.3.3), y esta ley es una ley de \emph{grupo
conmutativo} cuya simetría es $T(t)$, como resulta de la definición de un
objeto en grupos mediante diagramas conmutativos
($\mathbf 0_{\mathrm{III}}$, 8.2.5 y 8.2.6). Designaremos por
$\operatorname{Exan}_A(B,L)$ el grupo conmutativo así definido y diremos que
es el \emph{grupo de las clases de $A$-extensiones de $B$ por $L$}.
\end{env}

\begin{env}[18.3.5]
\label{0.18.3.5-es}
Designemos por $\mathbf K$ la categoría cuyos objetos son las ternas
$(A,B,L)$, donde $A$ es un anillo, $B$ un $A$-anillo y $L$ un $B$-bimódulo;
los \emph{morfismos} de esta categoría son las ternas $(u,v,w)$, donde
$u:A'\to A$ y $v:B'\to B$ son dos homomorfismos de anillos que hacen
conmutativo el cuadrado izquierdo del diagrama
\[
  \xymatrix@C=2pc@R=1.5pc{
    A\ar[r] & B & L\ar[d]^w\\
    A'\ar[u]^u\ar[r] & B'\ar[u]^v & L'
  }
\]
donde las flechas horizontales son los homomorfismos estructurales; por último,
$w:L\to L'$ es un homomorfismo de grupos conmutativos tal que se tiene
$w(v(b')z)=b'w(z)$ y

\oldpage[0\textsubscript{IV-1}]{61}

$w(zv(b'))=w(z)b'$ cualesquiera que sean $z\in L$ y $b'\in B'$ (en otros
términos, $w$ es un homomorfismo de $B'$-bimódulos cuando se dota a $L$ de la
estructura de $B'$-bimódulo definida por $v$). La composición de los morfismos
se define por
$(u',v',w')\circ(u,v,w)=(u\circ u',v\circ v',w'\circ w)$, lo que se justifica
inmediatamente. Nos proponemos demostrar que
\begin{equation}
\label{0.18.3.5.1-es}
\tag{18.3.5.1}
  (A,B,L)\longmapsto\operatorname{Exan}_A(B,L)
\end{equation}
es un \emph{funtor covariante} de la categoría $\mathbf K$ en la categoría
$\mathbf{Ab}$ de los grupos conmutativos. Se trata, pues, para toda terna
$(u,v,w)$ como la anterior, de definir un homomorfismo de grupos conmutativos
\[
  (u,v,w)_*:\operatorname{Exan}_A(B,L)
  \longrightarrow\operatorname{Exan}_{A'}(B',L').
\]
En virtud de la definición de los morfismos en $\mathbf K$, se puede escribir
\[
  (u,v,w)=(1_{A'},1_{B'},w)\circ(1_{A'},v,1_L)
  \circ(u,1_B,1_L)
\]
donde, en el primer factor, $L$ está dotado de su estructura de $B'$-bimódulo
definida por $v$; definiremos, por tanto, primero $(u,v,w)_*$ cuando dos de los
homomorfismos $u$, $v$, $w$ sean iguales a la identidad.
\end{env}

\begin{env}[18.3.6]
\label{0.18.3.6-es}
Definiremos primero $(1_A,1_B,w)_*$ como la aplicación
\begin{equation}
\label{0.18.3.6.1-es}
\tag{18.3.6.1}
  w_*:\operatorname{Exan}_A(B,L)\longrightarrow
  \operatorname{Exan}_A(B,L')
\end{equation}
denotada por $T(w)$ en (18.3.2); es inmediato comprobar que se trata de un
homomorfismo de grupos, pues esta propiedad se expresa mediante la
conmutatividad de diagramas que son las imágenes por $T$ de diagramas análogos
para $L$ y $L'$.

La aplicación $(1_A,v,1_L)_*$ es la aplicación
\begin{equation}
\label{0.18.3.6.2-es}
\tag{18.3.6.2}
  v^*:\operatorname{Exan}_A(B,L)\longrightarrow
  \operatorname{Exan}_A(B',L)
\end{equation}
definida del modo siguiente: si $E$ es una $A$-extensión de $B$ por $L$, se ha
visto que $E\times_BB'$ es una $A$-extensión de $B'$ por $L$ (18.2.5) y que,
si se reemplaza $E$ por una $A$-extensión $A$-equivalente $E'$,
$E'\times_BB'$ es $A$-equivalente a $E\times_BB'$; la imagen por $v^*$ de la
clase de $E$ es la clase de $E\times_BB'$. Se comprueba inmediatamente que, si
$w:L\to L'$ es un homomorfismo de $B$-bimódulos, el diagrama
\begin{equation}
\label{0.18.3.6.3-es}
\tag{18.3.6.3}
  \xymatrix@C=3pc@R=2.4pc{
    \operatorname{Exan}_A(B,L)\ar[r]^{v^*}\ar[d]_{w_*}
      & \operatorname{Exan}_A(B',L)\ar[d]^{w_*}\\
    \operatorname{Exan}_A(B,L')\ar[r]_{v^*}
      & \operatorname{Exan}_A(B',L')
  }
\end{equation}
es conmutativo, considerándose $L$ y $L'$ como $B'$-bimódulos por medio de
$v$ en la columna derecha. Reemplazando $L$ y $L'$, respectivamente, por
$L\times L$ y $L$, y $w$ por la adición $s$ en $L$, se concluye que $v^*$ es
efectivamente un \emph{homomorfismo de grupos}.
\end{env}

\oldpage[0\textsubscript{IV-1}]{62}

\begin{env}[18.3.6.4]
\label{0.18.3.6.4-es}
Finalmente, la aplicación $(u,1_B,1_L)_*$ es la aplicación
\[
  u^*:\operatorname{Exan}_A(B,L)\longrightarrow
  \operatorname{Exan}_{A'}(B,L)
\]
que se obtiene haciendo corresponder a una $A$-extensión $E$ de $B$ por $L$
el anillo $E$ considerado como $A'$-\emph{anillo} por medio de $u$ (18.1.1), que
es evidentemente una $A'$-extensión de $B$ por $L$, considerándose también $B$
como $A'$-anillo por medio de $u$~; es claro que una $A$-equivalencia es
también una $A'$-equivalencia, de donde resulta la aplicación (18.3.6.4), que, para todo
homomorfismo $w:L\to L'$ de $B$-bimódulos, hace también conmutativo el diagrama
\begin{equation}
\label{0.18.3.6.5-es}
\tag{18.3.6.5}
  \xymatrix@C=3pc@R=2.4pc{
    \operatorname{Exan}_A(B,L)\ar[r]^{u^*}\ar[d]_{w_*}
      & \operatorname{Exan}_{A'}(B,L)\ar[d]^{w_*}\\
    \operatorname{Exan}_A(B,L')\ar[r]_{u^*}
      & \operatorname{Exan}_{A'}(B,L')
  }
\end{equation}
de donde se concluye, como antes, que $u^*$ es un homomorfismo de grupos.

Dicho esto, se pone
$(u,v,w)_*=(1_{A'},1_{B'},w)_*\circ(1_{A'},v,1_L)_*
\circ(u,1_B,1_L)_*$ y se verifica fácilmente, en razón de la conmutatividad
de los diagramas (18.3.6.3) y (18.3.6.5), que se tiene
\[
  (u\circ u',v\circ v',w'\circ w)_*
  =(u',v',w')_*\circ(u,v,w)_*,
\]
lo que termina de demostrar que (18.3.4.1) es un funtor.

La existencia del homomorfismo de grupos (18.3.6.1) muestra, en particular,
que si $E$ es una $A$-extensión \emph{trivial} de $B$ por $L$, la extensión
$E\oplus_LL'$ de $B$ por $L'$ definida en (18.2.8) es también trivial, lo que
se verifica, por lo demás, sin dificultad de manera directa.

Por otra parte, para todo elemento $z$ del \emph{centro} $Z$ de $B$, la homotecia
$h_z:y\mapsto yz=zy$ es un endomorfismo del $B$-bimódulo $L$, por lo que $(h_z)_*$
es un endomorfismo del grupo conmutativo $\operatorname{Exan}_A(B,L)$, y por
funtorialidad, estos endomorfismos definen sobre
$\operatorname{Exan}_A(B,L)$ una estructura canónica de $Z$-\emph{módulo}.
\end{env}

\begin{env}[18.3.7]
\label{0.18.3.7-es}
Sean $A$, $A'$ dos anillos, $u:A'\to A$ un homomorfismo, $B$ un
$A$-anillo y $L$ un $B$-bimódulo. El \emph{núcleo} del homomorfismo de
grupos
\[
  u^*:\operatorname{Exan}_A(B,L)\longrightarrow
  \operatorname{Exan}_{A'}(B,L)
\]
está formado, por definición, por las clases de $A$-extensiones de $B$ por $L$ que
son $A'$-\emph{triviales} cuando se consideran como $A'$-extensiones por
medio de $u$. Este núcleo se denota por
$\operatorname{Exan}_{A/A'}(B,L)$ cuando ello no da lugar a confusión.

Si $\Lambda$ es un anillo, y si se tiene un diagrama conmutativo de
$\Lambda$-\emph{homomorfismos} de $\Lambda$-anillos
\begin{equation}
\label{0.18.3.7.1-es}
\tag{18.3.7.1}
  \xymatrix@C=2pc@R=1.5pc{
    B'\ar[r] & B\\
    A'\ar[u]\ar[r] & A\ar[u]
  }
\end{equation}

\oldpage[0\textsubscript{IV-1}]{63}

se deducen canónicamente homomorfismos
\begin{equation}
\label{0.18.3.7.2-es}
\tag{18.3.7.2}
  \operatorname{Exan}_{A/\Lambda}(B,L)\longrightarrow
  \operatorname{Exan}_{A'/\Lambda}(B,L)\longrightarrow
  \operatorname{Exan}_{A'/\Lambda}(B',L)
\end{equation}
que proceden de la conmutatividad del diagrama
\[
  \xymatrix@C=2.8pc@R=2.3pc{
    \operatorname{Exan}_A(B,L)\ar[r]\ar[d]
      & \operatorname{Exan}_{A'}(B,L)\ar[r]\ar[d]
      & \operatorname{Exan}_{A'}(B',L)\ar[d]\\
    \operatorname{Exan}_{\Lambda}(B,L)\ar[r]
      & \operatorname{Exan}_{\Lambda}(B,L)\ar[r]
      & \operatorname{Exan}_{\Lambda}(B',L)
  }
\]
donde las flechas se deducen de las de (18.3.7.1) por funtorialidad.
\end{env}

\begin{proposition}[18.3.8]
\label{0.18.3.8-es}
Sean $B$ un anillo, $\mathfrak J$ un ideal bilateral de $B$ y
$C=B/\mathfrak J$ el anillo cociente~; $\mathfrak J/\mathfrak J^2$ está entonces
canónicamente provisto de una estructura de $C$-bimódulo. Para todo $C$-bimódulo
$L$, sea $\operatorname{Hom}_C(\mathfrak J/\mathfrak J^2,L)$ el grupo
aditivo de los homomorfismos de $C$-bimódulos de $\mathfrak J/\mathfrak J^2$
en $L$. Se define entonces un isomorfismo canónico de grupos conmutativos
\begin{equation}
\label{0.18.3.8.1-es}
\tag{18.3.8.1}
  \eta_L:\operatorname{Hom}_C(\mathfrak J/\mathfrak J^2,L)
  \xrightarrow{\sim}\operatorname{Exan}_B(C,L)
\end{equation}
haciendo corresponder a todo $C$-homomorfismo
$w:\mathfrak J/\mathfrak J^2\to L$ (que es \emph{a fortiori} un
$B$-homomorfismo) la clase de la extensión
$(B/\mathfrak J^2)\oplus_{(\mathfrak J/\mathfrak J^2)}L$ deducida de
la extensión $B/\mathfrak J^2$ de $C$ por $\mathfrak J/\mathfrak J^2$ por medio
de $w$~; el isomorfismo inverso hace corresponder a la clase de una
$B$-extensión $E$ de $C$ por $L$ el homomorfismo $w$ tal que el compuesto
$\mathfrak J\to\mathfrak J/\mathfrak J^2\xrightarrow{w}L$ sea la restricción
a $\mathfrak J$ del homomorfismo estructural $B\to E$.
\end{proposition}

\begin{proof}
Sea $F=B/\mathfrak J^2$, que es una $B$-extensión de $C$ por
$\mathfrak J/\mathfrak J^2$. Para toda $B$-extensión
\[
  0\longrightarrow L\xrightarrow{j}E\xrightarrow{p}C\longrightarrow0
\]
de $C$ por $L$, el homomorfismo estructural $f:B\to E$ es tal que el compuesto
$B\xrightarrow{f}E\xrightarrow{p}C$ es el homomorfismo canónico
$B\to B/\mathfrak J=C$. Como la imagen de $\mathfrak J$ por $p\circ f$ es
nula, $f(\mathfrak J)$ está contenida en el núcleo de $p$, es decir,
$j(L)$, y como $j(L)$ es de cuadrado nulo, se tiene $f(\mathfrak J^2)=0$~; por tanto,
$f$ se factoriza como
$B\to B/\mathfrak J^2=F\xrightarrow{u}E$, y si $w$ es la restricción de
$u$ a $\mathfrak J/\mathfrak J^2$, se tiene un diagrama conmutativo
\[
  \xymatrix@C=2pc@R=1.5pc{
    0\ar[r] & \mathfrak J/\mathfrak J^2\ar[r]\ar[d]^w
      & F\ar[r]^g\ar[d]^u
      & C\ar[r]\ar[d]^{1_C}\ar@{}[d]|{\wr} & 0\\
    0\ar[r] & L\ar[r]_j & E\ar[r]_p & C\ar[r] & 0
  }
\]

dicho de otro modo, $(u,1_C,w)$ es un morfismo de extensiones (18.2.4). De él se
deduce un morfismo de extensiones
\[
  \xymatrix@C=2pc@R=1.5pc{
    0\ar[r] & L\ar[r]\ar[d]^{1_L}\ar@{}[d]|{\wr}
      & E'\ar[r]\ar[d]^{u'}
      & C\ar[r]\ar[d]^{1_C}\ar@{}[d]|{\wr} & 0\\
    0\ar[r] & L\ar[r]_j & E\ar[r]_p & C\ar[r] & 0
  }
\]

\oldpage[0\textsubscript{IV-1}]{64}

donde $E'=F\oplus_{(\mathfrak J/\mathfrak J^2)}L$ (18.2.8)~; se trata de
demostrar que $u'$ es una $B$-\emph{equivalencia}, lo que establecerá que $\eta_L$
es biyectivo. En virtud de la construcción realizada en (18.2.8), todo se reduce a
demostrar (habida cuenta de (18.1.7)) que la imagen inversa $G$ de la extensión
$E$ de $C$ por el homomorfismo canónico $g:F\to C$ es $F$-trivial, lo que
es evidente, puesto que $g$ se factoriza como
$F\xrightarrow{u}E\xrightarrow{p}C$ (18.1.6).

Queda ver que $\eta_L$ es un homomorfismo de \emph{grupos}~; ahora bien, para
todo $B$-homomorfismo $h:L\to L'$, es inmediato que el diagrama
\[
  \xymatrix@C=3pc@R=2.4pc{
    \operatorname{Hom}_C(\mathfrak J/\mathfrak J^2,L)
      \ar[r]^{\eta_L}\ar[d]_{\operatorname{Hom}(1,h)}
      & \operatorname{Exan}_B(C,L)\ar[d]^{h_*}\\
    \operatorname{Hom}_C(\mathfrak J/\mathfrak J^2,L')
      \ar[r]_{\eta_{L'}}^{\sim}
      & \operatorname{Exan}_B(C,L')
  }
\]
es conmutativo. Basta aplicar esta observación al homomorfismo
$L\times L\to L$ que define la adición para concluir.
\end{proof}

\subsection{Extensiones de álgebras}
\label{subsection:0.18.4-es}

\begin{env}[18.4.1]
\label{0.18.4.1-es}
Sea $A$ un anillo \emph{conmutativo}. La categoría de las $A$-\emph{álgebras} es
entonces una subcategoría plena de la de los $A$-anillos, caracterizada por el
hecho de que los homomorfismos estructurales $\rho:A\to B$ son \emph{centrales}.

Si $B$ es una $A$-álgebra y $L$ un $B$-bimódulo, es inmediato que la
$A$-extensión trivial tipo $D_B(L)$ es una $A$-álgebra. Del mismo modo, en la
construcción de (18.2.5) (resp. de (18.2.8)), si $B$, $B'$ y $E'$ (resp.
$B$ y $E$) son $A$-álgebras, ocurre lo mismo con
$E'\times_{B'}B$ (resp. con $E\oplus_LL'$). Finalmente, es claro que si $B$ es
una $A$-álgebra y $E$ una $A$-extensión de $B$ por un $B$-bimódulo $L$ que
es una $A$-álgebra, toda $A$-extensión de $B$ por $L$, $A$-\emph{equivalente}
a $E$, es también una $A$-álgebra. Se deduce entonces inmediatamente de la definición
(18.3.4) que las clases de $A$-extensiones equivalentes de $B$ por $L$ que
son $A$-\emph{álgebras} forman un \emph{subgrupo}, denotado por
$\operatorname{Exal}_A(B,L)$, de $\operatorname{Exan}_A(B,L)$. Sea
$\mathbf K'$ la subcategoría plena de la categoría $\mathbf K$ definida en
(18.3.5), cuyos objetos $(A,B,L)$ son tales que $A$ es conmutativo y $B$
una $A$-álgebra. Entonces lo que precede muestra que
\[
  (A,B,L)\longmapsto\operatorname{Exal}_A(B,L)
\]
es un \emph{funtor covariante} de $\mathbf K'$ en $\mathbf{Ab}$. Los
resultados de (18.3.7) no cambian cuando se reemplaza en todas partes
$\operatorname{Exan}$ por $\operatorname{Exal}$.
\end{env}

\begin{env}[18.4.2]
\label{0.18.4.2-es}
Supongamos todavía que el anillo $A$ es conmutativo. Las observaciones de (18.4.1) siguen
siendo válidas cuando se reemplaza «~álgebra~» por «~álgebra conmutativa~» y
«~bimódulo~» por «~módulo~». Si $B$ es una $A$-álgebra conmutativa y $L$ un
$B$-módulo,

\oldpage[0\textsubscript{IV-1}]{65}

las clases de $A$-extensiones equivalentes de $B$ por $L$ que son
$A$-\emph{álgebras conmutativas} (o, lo que equivale a lo mismo,
$A$-\emph{anillos conmutativos}) forman un \emph{subgrupo} de
$\operatorname{Exal}_A(B,L)$, denotado por $\operatorname{Exalcom}_A(B,L)$. Si
$\mathbf K''$ es la subcategoría plena de $\mathbf K'$ formada por los triples
$(A,B,L)$ donde $A$ y $B$ son conmutativos y $L$ un $B$-módulo,
\[
  (A,B,L)\longmapsto\operatorname{Exalcom}_A(B,L)
\]
es también un \emph{funtor covariante} de $\mathbf K''$ en $\mathbf{Ab}$. También se
puede, en (18.3.7), reemplazar en todas partes $\operatorname{Exan}$ por
$\operatorname{Exalcom}$. Finalmente, si en (18.3.8) se supone que $B$ es un
anillo \emph{conmutativo} y que $L$ es un $C$-\emph{módulo}, el mismo
razonamiento da un isomorfismo canónico
\begin{equation}
\label{0.18.4.2.1-es}
\tag{18.4.2.1}
  \operatorname{Hom}_C(\mathfrak J/\mathfrak J^2,L)
  \xrightarrow{\sim}\operatorname{Exalcom}_B(C,L)
\end{equation}
donde el primer miembro es el grupo de los homomorfismos de $C$-\emph{módulos}.
\end{env}

\begin{env}[18.4.3]
\label{0.18.4.3-es}
Sea $A$ un anillo conmutativo. Un caso importante de extensiones de
$A$-álgebras está formado por las $A$-álgebras $E$, extensiones de una $A$-álgebra $B$
por un $B$-bimódulo $L$ tales que $E$, como $A$-\emph{módulo}, es una
extensión \emph{trivial} del $A$-módulo $B$ por el $A$-módulo $L$~; dicho de otro
modo, la sucesión exacta de $A$-\emph{módulos}
$0\to L\to E\to B\to0$ es \emph{escindida}. Se dice entonces que $E$ es una
\emph{extensión de Hochschild} de $B$ por $L$. Este será siempre el caso cuando
$B$ sea un $A$-módulo \emph{proyectivo}, y en particular cuando $A$ sea un
\emph{cuerpo} conmutativo. Como $A$-módulo, se puede identificar $E$ con
$B\times L$, estando dada la multiplicación en $E$ por
$(x,s)(y,t)=(xy,xt+sy+f(x,y))$ con $f(x,y)\in L$. Si se expresa que esta
multiplicación define en $B\times L$ una estructura de $A$-álgebra, se obtiene
(M, XIV, 2) que $f$ debe ser una aplicación $A$-\emph{bilineal} de
$B\times B$ en $L$, tal que
\begin{equation}
\label{0.18.4.3.1-es}
\tag{18.4.3.1}
  f(xy,z)+f(x,y)z=xf(y,z)+f(x,yz),
\end{equation}
es decir, $f$ es un $2$-\emph{cociclo} sobre $B$ con valores en $L$, en el
sentido de Hochschild; para que la extensión $E$ sea $A$-trivial, es necesario
y suficiente que se tenga
\begin{equation}
\label{0.18.4.3.2-es}
\tag{18.4.3.2}
  f(x,y)=xg(y)-g(xy)+g(x)y,
\end{equation}
donde $g$ es una aplicación $A$-\emph{lineal} de $B$ en $L$; es decir, $f$
debe ser un $2$-\emph{coborde} en el sentido de Hochschild. Se deduce
inmediatamente que las clases de extensiones de Hochschild de $B$ por $L$
forman un subgrupo de $\operatorname{Exal}_A(B,L)$, isomorfo al \emph{grupo de
cohomología} de Hochschild $\mathrm H_A^2(B,L)$.

Si $B$ es una $A$-álgebra \emph{conmutativa}, y $L$ un $B$-módulo, las
extensiones de Hochschild \emph{conmutativas} de $B$ por $L$ corresponden a los
$2$-cociclos \emph{simétricos}, es decir, tales que $f(y,x)=f(x,y)$. Las clases
de extensiones de Hochschild conmutativas de $B$ por $L$ forman, por tanto, un
subgrupo de $\operatorname{Exalcom}_A(B,L)$, isomorfo al subgrupo de
$\mathrm H_A^2(B,L)$ que es imagen del grupo de los $2$-cociclos simétricos y
que denotaremos por $\mathrm H_A^2(B,L)^s$.
\end{env}

\begin{env}[18.4.4]
\label{0.18.4.4-es}
Limitémonos al caso en que $A$ y $B$ son \emph{conmutativos} y $L$ es un
$B$-módulo, y recordemos en este caso la definición equivalente de los grupos
de Hochschild $\mathrm H_A^i(B,L)$ (M, IX, 6). Se considera un complejo
$\mathbf P_\bullet=(\mathbf P_n)_{n\geqslant0}$ de $B$-módulos, donde
$\mathbf P_n=B^{\otimes(n+1)}=B\otimes_A B\otimes_A\cdots\otimes_A B$
($n+1$ veces), estando definida la estructura de $B$-módulo

\oldpage[0\textsubscript{IV-1}]{66}

por
\[
  x(y_1\otimes y_2\otimes\cdots\otimes y_{n+1})
  =(xy_1)\otimes y_2\otimes\cdots\otimes y_{n+1}\,;
\]
la derivación, de grado $-1$, $d_n:\mathbf P_n\to\mathbf P_{n-1}$, está
definida por
\begin{align*}
  d_n(x_1\otimes x_2\otimes\cdots\otimes x_{n+1})
    &=(x_1x_2)\otimes x_3\otimes\cdots\otimes x_{n+1}
      -x_1\otimes(x_2x_3)\otimes\cdots\otimes x_{n+1}+\cdots\\
    &\quad+(-1)^{n-1}x_1\otimes x_2\otimes\cdots\otimes x_{n-1}
      \otimes(x_nx_{n+1})
      +(-1)^n(x_{n+1}x_1)\otimes x_2\otimes\cdots\otimes x_n,
\end{align*}
que es efectivamente $B$-lineal puesto que $B$ es conmutativo.

Un $2$-cociclo de este complejo con valores en $L$ es una aplicación
$B$-lineal $h$ de $B\otimes B\otimes B$ en $L$ tal que
$h(d_3(x\otimes y\otimes z\otimes t))=0$; pero como
$h(x\otimes y\otimes z)=xh(1\otimes y\otimes z)$, el cociclo $h$ está
determinado por la aplicación $A$-bilineal
$(y,z)\longmapsto f(y,z)=h(1\otimes y\otimes z)$ de $B\times B$ en $L$, y,
escribiendo para $h$ la condición precedente, se recupera para $f$ la
condición (18.4.3.1). Del mismo modo, un $2$-coborde será una aplicación de la
forma $x\otimes y\otimes z\longmapsto h'(d_2(x\otimes y\otimes z))$, donde
$h':B\otimes_A B\to L$ es lineal, y aquí también $h'$ está determinado por la
aplicación $A$-lineal $g:y\longmapsto h'(1\otimes y)$ de $B$ en $L$; se
obtiene entonces
$h'(d_2(1\otimes x\otimes y))=xg(y)-g(xy)+yg(x)$, lo que vuelve a dar
(18.4.3.2). Se procede del mismo modo para todo $i$, y se ve así que se tiene
\begin{equation}
\label{0.18.4.4.1-es}
\tag{18.4.4.1}
  \mathrm H_A^i(B,L)=\mathrm H^i(\operatorname{Hom}_B(\mathbf P_\bullet,L)).
\end{equation}
\end{env}

\begin{env}[18.4.5]
\label{0.18.4.5-es}
Bajo las condiciones de (18.4.4), se puede interpretar del mismo modo el
grupo $\mathrm H_A^2(B,L)^s$ (18.4.3). Para ello, modifiquemos el complejo
$\mathbf P_\bullet$ en grado $3$, considerando un nuevo complejo
\[
  \mathbf P'_\bullet:\mathbf P'_3\xrightarrow{d'_3}\mathbf P_2
  \xrightarrow{d_2}\mathbf P_1\,;
\]
tomaremos
$\mathbf P'_3=\mathbf P_3\oplus(B\otimes_A B\otimes_A B)$, coincidiendo
$d'_3$ con $d_3$ sobre $\mathbf P_3$, y estando dado en
$B\otimes_A B\otimes_A B$ por
\[
  d'_3(x\otimes y\otimes z)=x\otimes y\otimes z-x\otimes z\otimes y.
\]

La relación $d_2d'_3=0$ resulta de la conmutatividad de $B$. Con las notaciones
introducidas anteriormente, un $2$-cociclo de $\mathbf P'_\bullet$ corresponde
ahora a una aplicación $A$-bilineal $f:B\times B\to L$ que es
\emph{simétrica} y verifica (18.4.3.1); por consiguiente, se tiene
\[
  \mathrm H_A^2(B,L)^s
  =\mathrm H^2(\operatorname{Hom}_B(\mathbf P'_\bullet,L)).
\]
\end{env}

\begin{env}[18.4.6]
\label{0.18.4.6-es}
En el caso particular en que se considera un \emph{cuerpo} conmutativo $k$ y
una \emph{extensión} $K$ de $k$, se tiene
$\operatorname{Ext}_K^1(M,L)=0$ cualesquiera que sean los $K$-espacios
vectoriales $L$, $M$, y, por consiguiente (M, VI, 3.3 $a$)), se tiene un
isomorfismo canónico
\begin{equation}
\label{0.18.4.6.1-es}
\tag{18.4.6.1}
  \mathrm H_k^i(K,L)\xrightarrow{\sim}
  \operatorname{Hom}_K(\mathrm H_i(\mathbf P_\bullet),L)
\end{equation}
y, del mismo modo,
\begin{equation}
\label{0.18.4.6.2-es}
\tag{18.4.6.2}
  \mathrm H_k^2(K,L)^s\xrightarrow{\sim}
  \operatorname{Hom}_K(\mathrm H_2(\mathbf P'_\bullet),L).
\end{equation}
\end{env}

\subsection{Caso de los anillos topológicos}
\label{subsection:0.18.5-es}

\begin{env}[18.5.1]
\label{0.18.5.1-es}
Sean $A$, $B$ dos anillos \emph{topológicos} cuya topología es \emph{lineal},
$\rho:A\to B$ un homomorfismo continuo y $L$ un $B$-bimódulo topológico, y
supongamos que existe un ideal bilateral \emph{abierto} $\mathfrak R_0$ de $B$
tal que $\mathfrak R_0.L=L.\mathfrak R_0=0$, de modo que $L$ puede ser

\oldpage[0\textsubscript{IV-1}]{67}

considerado como un $(B/\mathfrak R)$-bimódulo para todo ideal abierto
$\mathfrak R\subset\mathfrak R_0$. Sean entonces
$\mathfrak R\subset\mathfrak R_0$ un ideal bilateral abierto de $B$,
$\mathfrak J$ un ideal bilateral abierto de $A$ tal que
$\rho(\mathfrak J)\subset\mathfrak R$, de modo que $B/\mathfrak R$ puede ser
considerado como un $(A/\mathfrak J)$-anillo discreto. La observación
precedente prueba que el grupo
$\operatorname{Exan}_{A/\mathfrak J}(B/\mathfrak R,L)$ está definido; además,
si $\mathfrak R'\subset\mathfrak R$, $\mathfrak J'\subset\mathfrak J$ son dos
ideales bilaterales abiertos de $B$ y $A$, respectivamente, tales que
$\rho(\mathfrak J')\subset\mathfrak R'$, se tiene por (18.3.5.1) un
homomorfismo canónico
\begin{equation}
\label{0.18.5.1.1-es}
\tag{18.5.1.1}
  \operatorname{Exan}_{A/\mathfrak J}(B/\mathfrak R,L)
  \longrightarrow
  \operatorname{Exan}_{A/\mathfrak J'}(B/\mathfrak R',L).
\end{equation}

El conjunto de pares de ideales abiertos $(\mathfrak J,\mathfrak R)$ tales que
$\rho(\mathfrak J)\subset\mathfrak R\subset\mathfrak R_0$ está evidentemente
ordenado y es filtrante por la derecha para la relación «~$\mathfrak J\supset\mathfrak J'$
y $\mathfrak R\supset\mathfrak R'$~», y las aplicaciones (18.5.1.1)
definen un \emph{sistema inductivo} de grupos aditivos cuyo conjunto
de índices es este conjunto. Se pone, por abuso de notación (pues ya no se trata
de un grupo en correspondencia biunívoca natural con un conjunto
de extensiones),
\begin{equation}
\label{0.18.5.1.2-es}
\tag{18.5.1.2}
  \operatorname{Exantop}_A(B,L)
  =\varinjlim\operatorname{Exan}_{A/\mathfrak J}(B/\mathfrak R,L).
\end{equation}

Decir que el segundo miembro de (18.5.1.2) es \emph{nulo} significa, pues, que, para
todo par de ideales abiertos $\mathfrak J\subset A$, $\mathfrak R\subset B$
tales que $\rho(\mathfrak J)\subset\mathfrak R\subset\mathfrak R_0$ y toda
$(A/\mathfrak J)$-extensión $E$ de $B/\mathfrak R$ por $L$, existen dos
ideales abiertos $\mathfrak J'\subset\mathfrak J$,
$\mathfrak R'\subset\mathfrak R$ tales que
$\rho(\mathfrak J')\subset\mathfrak R'$ y que la imagen inversa por
el homomorfismo $B/\mathfrak R'\to B/\mathfrak R$ de $E$ sea \emph{trivial}.

Dejamos al lector la definición análoga de los límites inductivos
$\operatorname{Exaltop}_A(B,L)$, $\operatorname{Exalcotop}_A(B,L)$ a partir de
$\operatorname{Exal}_{A/\mathfrak J}(B/\mathfrak R,L)$ y de
$\operatorname{Exalcom}_{A/\mathfrak J}(B/\mathfrak R,L)$ para el caso en que $A$
es conmutativo y $B$ una $A$-álgebra (resp. una $A$-álgebra conmutativa)
topológica.
\end{env}

\begin{env}[18.5.2]
\label{0.18.5.2-es}
Si se tiene un diagrama conmutativo de homomorfismos continuos de anillos
\[
  \xymatrix@C=3pc@R=2pc{
    B'\ar[r] & B\\
    A'\ar[u]\ar[r] & A\ar[u]
  }
\]
se deducen canónicamente dos homomorfismos de grupos aditivos
\[
  \operatorname{Exantop}_A(B,L)
  \longrightarrow\operatorname{Exantop}_{A'}(B,L)
  \longrightarrow\operatorname{Exantop}_{A'}(B',L)
\]
pasando al límite inductivo a partir de (18.3.5.1).

En virtud de la exactitud del funtor $\varinjlim$ en la categoría de los grupos
conmutativos, el núcleo del homomorfismo
\[
  \operatorname{Exantop}_A(B,L)
  \longrightarrow\operatorname{Exantop}_{A'}(B,L)
\]
es el límite inductivo de los núcleos de los homomorfismos
\[
  \operatorname{Exan}_{A/\mathfrak J}(B/\mathfrak R,L)
  \longrightarrow
  \operatorname{Exan}_{A'/\mathfrak J'}(B/\mathfrak R,L)
\]
donde se ha tomado como $\mathfrak J'$ la imagen inversa de $\mathfrak J$; se denota
este núcleo por $\operatorname{Exantop}_{A/A'}(B,L)$. Se definen del mismo modo
$\operatorname{Exaltop}_{A/A'}(B,L)$ y
$\operatorname{Exalcotop}_{A/A'}(B,L)$. Finalmente, si se tiene un homomorfismo

\oldpage[0\textsubscript{IV-1}]{68}

continuo de $B$-bimódulos $L\to L'$, se deduce canónicamente un
homomorfismo de grupos aditivos
\[
  \operatorname{Exantop}_A(B,L)\longrightarrow
  \operatorname{Exantop}_A(B,L')
\]
pasando al límite inductivo a partir de (18.3.6.1).
\end{env}

\begin{env}[18.5.3]
\label{0.18.5.3-es}
Dados un anillo topológico $C$ y dos $C$-bimódulos topológicos
$M$, $N$,\linebreak se denota por $\operatorname{Hom.\ cont}_C(M,N)$ el grupo aditivo de los
$C$-homomorfismos \emph{continuos} de $M$ en $N$.
\end{env}

\begin{lemma}[18.5.3.1]
\label{0.18.5.3.1-es}
Sean $C$ un anillo topológico y $E$, $L$ dos $C$-bimódulos topológicos;
se supone que las topologías son lineales, que $L$ es discreto y está anulado
por un ideal bilateral abierto de $C$. Entonces se tiene un isomorfismo canónico
\begin{equation}
\label{0.18.5.3.2-es}
\tag{18.5.3.2}
  \varinjlim\operatorname{Hom}_{C/\mathfrak R}(E/V,L)
  \xrightarrow{\sim}\operatorname{Hom.\ cont}_C(E,L)
\end{equation}
donde en el primer miembro el límite inductivo se toma con respecto al conjunto
ordenado filtrante de pares $(\mathfrak R,V)$ tales que $\mathfrak R$ es un
ideal bilateral abierto de $C$, $V$ un sub-$C$-bimódulo abierto de $E$, tales que
$\mathfrak R.L=L.\mathfrak R=0$, $\mathfrak R.E\subset V$,
$E.\mathfrak R\subset V$.
\end{lemma}

\begin{proof}
Como $C/\mathfrak R$ y $E/V$ son discretos, se tienen homomorfismos
canónicos
\[
  w_{\mathfrak R,V}:\operatorname{Hom}_{C/\mathfrak R}(E/V,L)
  \longrightarrow\operatorname{Hom.\ cont}_C(E,L)
\]
que forman un sistema inductivo, de donde se obtiene un homomorfismo (18.5.3.2) pasando al
límite inductivo. Como el homomorfismo $E/V'\to E/V$ es sobreyectivo para
$V'\subset V$, de la definición resulta inmediatamente que el homomorfismo
\[
  \operatorname{Hom}_{C/\mathfrak R}(E/V,L)
  \longrightarrow\operatorname{Hom}_{C/\mathfrak R}(E/V',L)
\]
(con $\mathfrak R.L=L.\mathfrak R=0$, $\mathfrak R.E\subset V'$,
$E.\mathfrak R\subset V'$) es inyectivo, y evidentemente sucede lo mismo con
el homomorfismo
\[
  \operatorname{Hom}_{C/\mathfrak R}(E/V,L)
  \longrightarrow\operatorname{Hom}_{C/\mathfrak R'}(E/V,L)
\]
para $\mathfrak R'\subset\mathfrak R$; se concluye que el homomorfismo
(18.5.3.2) es inyectivo. Por otra parte, si $u$ es un $C$-homomorfismo continuo
de $E$ en $L$, su núcleo es un subbimódulo abierto $V_0$ de $E$, y si
$\mathfrak R_0$ es un ideal bilateral abierto de $C$ tal que
$\mathfrak R_0.L=L.\mathfrak R_0=0$ y $\mathfrak R_0.E\subset V_0$,
$E.\mathfrak R_0\subset V_0$, está claro que $u$ es la imagen canónica de un
$(C/\mathfrak R_0)$-homomorfismo de $E/V_0$ en $L$, luego (18.5.3.2) es
sobreyectivo.
\end{proof}

Dicho esto, la proposición (18.3.8) se generaliza como sigue a los anillos
topológicos:

\begin{proposition}[18.5.4]
\label{0.18.5.4-es}
Sean $B$ un anillo topológico linealmente topologizado, $\mathfrak J$ un
ideal bilateral de $B$ y $C=B/\mathfrak J$ el anillo topológico cociente;
$\mathfrak J/\mathfrak J^2$ (donde $\mathfrak J$ está dotado de la topología inducida
por la de $B$ y $\mathfrak J/\mathfrak J^2$ de la topología cociente de
la de $\mathfrak J$) queda entonces canónicamente dotado de una estructura de
$C$-bimódulo topológico. Para todo $C$-bimódulo discreto $L$ anulado por un
ideal abierto de $C$, existe entonces un isomorfismo canónico
\begin{equation}
\label{0.18.5.4.1-es}
\tag{18.5.4.1}
  \operatorname{Hom.\ cont}_C(\mathfrak J/\mathfrak J^2,L)
  \xrightarrow{\sim}\operatorname{Exantop}_B(C,L).
\end{equation}
\end{proposition}

\begin{proof}
En efecto, para todo ideal abierto $\mathfrak R$ de $B$ tal que
$(\mathfrak J+\mathfrak R)/\mathfrak J$ anule $L$, se tiene, según (18.3.9),
un isomorfismo canónico
\[
  \operatorname{Hom}_{B/(\mathfrak J+\mathfrak R)}
    ((\mathfrak J+\mathfrak R)/(\mathfrak J^2+\mathfrak R),L)
  \xrightarrow{\sim}
  \operatorname{Exan}_{B/\mathfrak R}(B/(\mathfrak J+\mathfrak R),L)
\]
y basta pasar al límite inductivo, teniendo en cuenta (18.5.3.1).
\end{proof}
